\chapter{Literature Review}
\label{ch:literature-review}

This chapter surveys prior work relevant to four design pillars of CAAS:
(i) air-quality forecasting for fine particulate matter, (ii) model-family
choice for tabular time-series problems, (iii) the role of satellite fire
observations in atmospheric prediction, and (iv) production-grade MLOps
patterns for environmental machine-learning systems. Each section motivates
a concrete choice in the CAAS design.

\section{Air-Quality Forecasting for PM2.5}

Fine particulate matter with aerodynamic diameter below 2.5~\textmu m
(PM2.5) is the single largest environmental cause of premature mortality
worldwide, with the World Health Organization (\cite{who2021}) revising its
24-hour guideline downward to 15~\textmu g/m\textsuperscript{3} in 2021 to
reflect strengthened epidemiological evidence. Government monitoring networks
in Thailand publish measured PM2.5 through portals such as air4thai
(\cite{pcd}), but these portals are retrospective --- they report what was
observed, not what is expected over the coming days. This gap motivates the
present work: forecast-driven alerting.

Statistical and machine-learning approaches to PM2.5 forecasting have
evolved from simple persistence and ARIMA baselines toward gradient-boosted
tree ensembles and deep sequence models. For Southeast Asian haze events,
\cite{chaing_mai_haze} demonstrated that regional meteorological and fire
drivers dominate PM2.5 dynamics during the burning season, motivating
multi-source feature fusion. Satellite-derived proxies such as aerosol
optical depth have also been calibrated against ground PM2.5 measurements
(\cite{psi}), though such proxies are coarse in time and are not directly
used in CAAS because the Chiang Mai monitoring station provides a direct
PM2.5 measurement at daily resolution.

CAAS adopts a daily tabular formulation with engineered lag, rolling, and
cyclic-calendar features, deferring image-based aerosol retrievals to future
work. The persistence baseline (yesterday's value) is retained as a
reference against which any learned model must demonstrate a meaningful
improvement.

\section{Gradient Boosting Versus Deep Learning for Tabular Time Series}

The question of whether to use gradient-boosted decision trees or deep
sequence models on tabular time-series problems remains contested in the
applied literature. XGBoost (\cite{xgboost}) and LightGBM --- the two
dominant gradient-boosting implementations --- offer histogram-based
training, native handling of missing values, and strong regularisation
controls, with inference latencies typically in the low milliseconds for
tables of this size. Long Short-Term Memory (LSTM) networks
(\cite{lstm}) in contrast explicitly model temporal dependencies through
gated recurrence, at the cost of greater hyperparameter sensitivity and
higher training budgets.

Recent empirical surveys have repeatedly found that on moderately sized
tabular problems with engineered temporal features, gradient-boosted models
match or outperform deep sequence models. The value of deep models tends to
appear only at very long sequence lengths, in multi-variate settings with
complex cross-feature interactions, or when raw un-engineered features
dominate. The CAAS feature matrix (45 engineered features, $\sim$5,000 daily
records) sits squarely within the regime where gradient-boosted methods
typically win.

This literature informs the CAAS model-selection hierarchy: LightGBM is
designated champion, XGBoost is retained as a strong-secondary for runtime
A/B comparison, and LSTM is trained as a comparison baseline to test --- not
assume --- whether temporal sequence structure adds value over
tree-friendly lag features. Chapter~\ref{ch:ml-development} reports the
empirical outcome of this comparison.

\section{Satellite Fire Observations in Atmospheric Modelling}

The NASA Fire Information for Resource Management System (FIRMS,
\cite{firms}) delivers near-real-time active-fire detections from VIIRS and
MODIS sensors, including fire radiative power (FRP) as an intensity proxy.
FIRMS data have been used extensively in smoke-transport modelling and
regional air-quality forecasting, particularly for agricultural-burning
events in Southeast Asia where in-situ fire monitoring is sparse.

For Chiang Mai specifically, the dominant PM2.5 source during February--April
is biomass burning in the surrounding highlands and in neighbouring countries
upwind. Incorporating FIRMS hotspot counts and rolling fire accumulation
features is therefore physically motivated: yesterday's fire activity
influences tomorrow's PM2.5 concentration through transport and atmospheric
mixing. The CAAS feature engineering reflects this mechanism through
hotspot counts within 50~km and 100~km radii and 7-day/14-day rolling
windows. Chapter~\ref{ch:ml-development} reports an ablation study
(Scenario~D) that quantifies the measurable contribution of these
FIRMS-derived features relative to the weather-and-lag-only baseline.

\section{MLOps Patterns for Environmental Machine-Learning Systems}

Operationalising a machine-learning model involves far more than training.
The MLOps stack adopted in CAAS follows the now-standard pattern in which
distinct tools handle distinct concerns: MLflow (\cite{mlflow}) for
experiment tracking and model registry; FastAPI (\cite{fastapi}) for
low-latency model serving; Evidently (\cite{evidently}) for distribution
drift monitoring; Terraform (\cite{terraform}) for cloud infrastructure
as code; Docker (\cite{docker}) for reproducible containerised deployment;
and GitHub Actions (\cite{githubactions}) for CI/CD orchestration.

Drift detection in environmental ML systems presents a specific complication
that generic MLOps literature does not address: strong, expected seasonal
variation. A naive population-stability-index (PSI) or
Kolmogorov--Smirnov (\cite{ks_test}) drift alarm raised against a summer
baseline will fire during the January--April haze season every year, producing
false positives and eroding operator trust. CAAS addresses this with a
two-tier drift policy that combines core-feature thresholds with seasonal
reference windows, described in Chapter~\ref{ch:mlops}. This seasonal-aware
retraining policy is, to our knowledge, a novel contribution in the context
of student-scale environmental MLOps projects and is one of the key
engineering contributions of this work.
